\section*{Лекция 5. Введение в задачу классификации}
\subsection*{Введение}

Итак, ранее мы уже вводили определение отступа в задаче классификации. Напомним, что 
\[
M = y_i \langle w, x \rangle
\]

и эта величина характеризует неотнормированное расстояние от точки до гиперплоскости. Далее мы вводили 
\[
L = [M < 0]
\]
как функцию потерь. В силу невозможности оптимизации для этой функции потерь, мы предположили, что было бы неплохо вводить т.н. \textit{оценку} данной функции потерь.

На прошлой лекции мы привели пару примеров возможных оценок, и сегодня мы остановимся на 
\[
\tilde{M} = \log(1 + \exp^{-M})
\]

На сегодняшнем занятии мы постараемся объяснить алгоритм, который стоит за данной оценкой, и почему в действительности данная оценка имеет место быть.